% ============================================================
%  BIPOL 2 — Szablon Projektu
%  Lotnictwo / Mechanika Lotu
% ============================================================
\documentclass[12pt, a4paper]{article}

% ── Pakiety ─────────────────────────────────────────────────
\usepackage[T1]{fontenc}
\usepackage[utf8]{inputenc}
\usepackage[polish]{babel}
\usepackage{geometry}
\usepackage{amsmath, amssymb}
\usepackage{graphicx}
\usepackage{booktabs}
\usepackage{array}
\usepackage{tabularx}
\usepackage{float}
\usepackage{caption}
\usepackage{subcaption}
\usepackage{hyperref}
\usepackage{setspace}
\usepackage{fancyhdr}
\usepackage{titlesec}
\usepackage{enumitem}
\usepackage{xcolor}
\usepackage{soul}          % podświetlenie \hl{}
\usepackage{siunitx}       % jednostki SI
\usepackage{multirow}
\usepackage{longtable}

\geometry{
  top    = 2.5cm,
  bottom = 2.5cm,
  left   = 2.5cm,
  right  = 2.5cm
}

% ── Nagłówki sekcji ─────────────────────────────────────────
\titleformat{\section}{\large\bfseries}{\thesection.}{0.5em}{}
\titleformat{\subsection}{\normalsize\bfseries}{\thesubsection.}{0.5em}{}
\titleformat{\subsubsection}{\normalsize\itshape}{\thesubsubsection.}{0.5em}{}

% ── Nagłówek / stopka ────────────────────────────────────────
\pagestyle{fancy}
\fancyhf{}
\fancyhead[L]{\small BIPOL 2 — \tytulprojektu}
\fancyhead[R]{\small \imienazwisko}
\fancyfoot[C]{\thepage}
\renewcommand{\headrulewidth}{0.4pt}

% ── Dane dokumentu ───────────────────────────────────────────
%  Wypełnij poniższe pola przed kompilacją
\newcommand{\imienazwisko}{Imię Nazwisko}
\newcommand{\nrindeksu}{XXXXXX}
\newcommand{\grupa}{X — dzień HH:MM–HH:MM}
\newcommand{\prowadzacy}{Tytuł Nazwisko}
\newcommand{\nrprojektu}{X}
\newcommand{\tytulprojektu}{Tytuł projektu}
\newcommand{\dataoddania}{\hl{XX.03.2026}}

% ── Hiperłącza ───────────────────────────────────────────────
\hypersetup{
  colorlinks = true,
  linkcolor  = black,
  urlcolor   = blue,
  citecolor  = black
}

% ============================================================
\begin{document}

% ── Strona tytułowa ─────────────────────────────────────────
\begin{titlepage}
  \centering
  \vspace*{1.5cm}

  {\LARGE\bfseries BIPOL 2 \par}
  \vspace{0.4cm}
  {\large Podstawy i Budowa Obiektów Latających 2 \par}
  \vspace{1.5cm}

  {\Large\bfseries Projekt \nrprojektu: \par}
  \vspace{0.4cm}
  {\Large \tytulprojektu \par}
  \vspace{2.5cm}

  \begin{tabular}{rl}
    \textbf{Autor:}        & \imienazwisko \\[4pt]
    \textbf{Nr indeksu:}   & \nrindeksu    \\[4pt]
    \textbf{Grupa:}        & \grupa        \\[4pt]
    \textbf{Prowadzący:}   & \prowadzacy   \\[10pt]
    \textbf{Data oddania:} & \dataoddania  \\[10pt]
    \textbf{Ocena:}        & \dotfill      \\
  \end{tabular}

  \vfill
  \today
\end{titlepage}

% ── Spis treści ─────────────────────────────────────────────
\tableofcontents
\newpage

% ============================================================
%  1. Dane geometryczne usterzenia
% ============================================================
\section{Dane geometryczne usterzenia}

Na podstawie podobnych konstrukcji samolotów dyspozycyjnych określono wstępne wymiary
usterzenia (Projekt~3 z BIPOL~1).

\subsection{Usterzenie poziome}

% Tabela przykładowa — uzupełnij wartości
\begin{table}[H]
  \centering
  \caption{Podstawowe parametry geometryczne usterzenia poziomego}
  \label{tab:geom_poziome}
  \begin{tabular}{lcc}
    \toprule
    \textbf{Parametr} & \textbf{Symbol} & \textbf{Wartość} \\
    \midrule
    Powierzchnia usterzenia & $S_H$          & \SI{}{m^2}      \\
    Rozpiętość              & $b_H$          & \SI{}{m}        \\
    Wydłużenie              & $\Lambda_H$    & ---             \\
    Skos krawędzi natarcia  & $\varphi_{H}$  & \SI{}{\degree}  \\
    Ramię usterzenia        & $l_H$          & \SI{}{m}        \\
    Współczynnik objętości  & $V_H$          & ---             \\
    \bottomrule
  \end{tabular}
\end{table}

\subsection{Stery wysokości}

% Uzupełnij dane steru
\begin{table}[H]
  \centering
  \caption{Parametry steru wysokości}
  \label{tab:ster}
  \begin{tabular}{lcc}
    \toprule
    \textbf{Parametr} & \textbf{Symbol} & \textbf{Wartość} \\
    \midrule
    Cięciwa steru (śr.) & $c_e$   & \SI{}{m}  \\
    Stosunek cięciw     & $c_e/c$ & ---       \\
    Kąt wychylenia (+)  & $\delta_{e,\max}^+$ & \SI{}{\degree} \\
    Kąt wychylenia (-)  & $\delta_{e,\max}^-$ & \SI{}{\degree} \\
    \bottomrule
  \end{tabular}
\end{table}

% ============================================================
%  2. Moment pochylający samolotu
% ============================================================
\section{Moment pochylający samolotu}

Określono moment pochylający samolotu. Wyznaczono obwiednie obciążeń dopuszczalnych
oraz obwiednie obciążeń od mechanizacji.

\subsection{Współczynnik momentu pochylającego}

\begin{equation}
  C_m = C_{m_0} + C_{m_\alpha}\,\alpha + C_{m_{\delta_e}}\,\delta_e
  \label{eq:Cm}
\end{equation}

% ── Opis symboli ──
\noindent
gdzie:
\begin{itemize}[noitemsep]
  \item $C_{m_0}$ — współczynnik momentu przy $\alpha = 0$,
  \item $C_{m_\alpha}$ — pochodna momentu pochylającego po kącie natarcia \SI{}{1/\radian},
  \item $C_{m_{\delta_e}}$ — pochodna momentu pochylającego po wychyleniu steru \SI{}{1/\radian}.
\end{itemize}

\subsection{Wyniki}

% Wstaw wykres lub tabelę z wynikami
\begin{figure}[H]
  \centering
  \includegraphics[width=0.8\textwidth]{figures/moment_pochylajacy.pdf}
  \caption{Moment pochylający samolotu w funkcji kąta natarcia}
  \label{fig:Cm}
\end{figure}

% ============================================================
%  3. Charakterystyki aerodynamiczne usterzenia poziomego
% ============================================================
\section{Charakterystyki aerodynamiczne usterzenia poziomego}

Wyznaczono obwiednie obciążeń dopuszczalnych oraz obwiednie obciążeń od mechanizacji.

\subsection{Gradient siły nośnej}

\begin{equation}
  C_{L_H} = C_{L_{H,\alpha_H}} \left( \alpha - \varepsilon - i_H + \frac{\delta_e\,C_{L_{H,\delta_e}}}{C_{L_{H,\alpha_H}}} \right)
  \label{eq:CLH}
\end{equation}

% ============================================================
%  4. Punkty neutralne stateczności
% ============================================================
\section{Punkty neutralne stateczności}

Wyznaczono obwiednie obciążeń dopuszczalnych oraz obwiednie obciążeń od mechanizacji.

\subsection{Punkt neutralny przy sterze swobodnym}

\begin{equation}
  \bar{x}_{n,\text{free}} = \bar{x}_{n,\text{fixed}} - \frac{C_{H_\alpha}}{C_{H_{\delta_e}}} \cdot
  \frac{C_{L_{H,\delta_e}}}{C_{L_\alpha}} \cdot \frac{S_H l_H}{S l}
  \label{eq:xn_free}
\end{equation}

\subsection{Punkt neutralny przy sterze ustalonym}

% Uzupełnij wyprowadzenie lub tabelę wyników
\begin{table}[H]
  \centering
  \caption{Punkty neutralne stateczności}
  \label{tab:pkt_neutralne_stat}
  \begin{tabular}{lcc}
    \toprule
    \textbf{Przypadek} & \textbf{Symbol} & $\bar{x}_n$ \\
    \midrule
    Ster ustalony & $\bar{x}_{n,\text{fixed}}$ & \\
    Ster swobodny & $\bar{x}_{n,\text{free}}$  & \\
    \bottomrule
  \end{tabular}
\end{table}

% ============================================================
%  5. Punkty neutralne sterowności
% ============================================================
\section{Punkty neutralne sterowności}

Wyznaczono obwiednie obciążeń dopuszczalnych oraz obwiednie obciążeń od mechanizacji.

\subsection{Warunek sterowności}

\begin{equation}
  \frac{\partial F_s}{\partial n} = 0 \quad \Rightarrow \quad
  \bar{x}_{nm} = \ldots
  \label{eq:sterownosc}
\end{equation}

% ============================================================
%  6. Ocena wyważenia — gradienty sił i wychyleń steru wysokości
% ============================================================
\section{Ocena wyważenia samolotu — gradienty sił i wychyleń steru wysokości}

Wyznaczono obwiednie obciążeń dopuszczalnych oraz obwiednie obciążeń od mechanizacji.

\subsection{Gradient wychylenia steru}

\begin{equation}
  \left.\frac{d\delta_e}{dn}\right|_{F_s=\text{const}} = \ldots
\end{equation}

\subsection{Gradient siły na drążku}

\begin{equation}
  \left.\frac{dF_s}{dn}\right|_{\delta_e=\text{const}} = \ldots
\end{equation}

\begin{table}[H]
  \centering
  \caption{Gradienty dla różnych konfiguracji lotu}
  \label{tab:gradienty}
  \begin{tabularx}{\textwidth}{Xccc}
    \toprule
    \textbf{Konfiguracja} & $d\delta_e/dn$ [\SI{}{\degree}] &
    $dF_s/dn$ [\si{N}] & \textbf{Ocena} \\
    \midrule
    Podejście (klapy wychylone)  & & & \\
    Lot ustalony (klapy schowane) & & & \\
    \bottomrule
  \end{tabularx}
\end{table}

% ============================================================
%  7. Równowaga samolotu
% ============================================================
\section{Równowaga samolotu}

Wyznaczono obwiednie obciążeń dopuszczalnych oraz obwiednie obciążeń od mechanizacji.

\subsection{Warunki równowagi podłużnej}

\begin{align}
  C_L &= C_{L,\alpha}\,\alpha + C_{L,\delta_e}\,\delta_e + C_{L_0} \label{eq:CL_rownow} \\
  C_m &= 0 \label{eq:Cm_rownow}
\end{align}

\subsection{Wykresy równowagi}

\begin{figure}[H]
  \centering
  \begin{subfigure}[b]{0.48\textwidth}
    \includegraphics[width=\textwidth]{figures/rownowaga_delta.pdf}
    \caption{Wychylenie steru w funkcji $C_L$}
  \end{subfigure}
  \hfill
  \begin{subfigure}[b]{0.48\textwidth}
    \includegraphics[width=\textwidth]{figures/rownowaga_sila.pdf}
    \caption{Siła na drążku w funkcji $C_L$}
  \end{subfigure}
  \caption{Charakterystyki równowagi samolotu}
  \label{fig:rownowaga}
\end{figure}

% ============================================================
%  8. Podsumowanie
% ============================================================
\section{Podsumowanie}

Wyznaczono obwiednie obciążeń dopuszczalnych oraz obwiednie obciążeń od mechanizacji.

\begin{itemize}
  \item Stateczność podłużna jest zapewniona dla wszystkich badanych konfiguracji, ponieważ
        punkt neutralny znajduje się za środkiem ciężkości ($\bar{x}_n > \bar{x}_{cg}$).
  \item Gradienty wychylenia steru i sił na drążku mieszczą się w wymaganych zakresach.
  \item \textit{[Uzupełnij pozostałe wnioski]}
\end{itemize}

% ============================================================
%  Literatura
% ============================================================
\newpage
\bibliographystyle{plain}
% \bibliography{literatura}  % odkomentuj i utwórz plik literatura.bib

\begin{thebibliography}{9}
  \bibitem{raymer} Raymer, D.P., \textit{Aircraft Design: A Conceptual Approach}, AIAA, 2018.
  \bibitem{etkin} Etkin, B., Reid, L.D., \textit{Dynamics of Flight: Stability and Control}, Wiley, 1996.
  \bibitem{materialy} Materiały dydaktyczne BIPOL~2, Politechnika Warszawska.
\end{thebibliography}

% ============================================================
%  Załączniki (opcjonalnie)
% ============================================================
\newpage
\appendix
\section{Kod obliczeniowy / skrypty}

% Wstaw kod lub odniesienie do pliku

\section{Dane wejściowe — zestawienie}

\begin{table}[H]
  \centering
  \caption{Kompletne dane wejściowe projektu}
  \label{tab:dane_wejsciowe}
  \begin{tabular}{llcc}
    \toprule
    \textbf{Parametr} & \textbf{Opis} & \textbf{Symbol} & \textbf{Wartość} \\
    \midrule
    Masa startowa        & MTOW                    & $m_{max}$    & \SI{}{kg}   \\
    Powierzchnia skrzydła & Wing area               & $S$          & \SI{}{m^2}  \\
    Rozpiętość skrzydła  & Wing span               & $b$          & \SI{}{m}    \\
    Środkowa cięciwa     & MAC                     & $\bar{c}$    & \SI{}{m}    \\
    Prędkość przelotowa  & Cruise speed            & $V_{cr}$     & \SI{}{m/s}  \\
    Położenie CG (przód) & Forward CG              & $\bar{x}_{cg,f}$ & ---    \\
    Położenie CG (tył)   & Aft CG                  & $\bar{x}_{cg,a}$ & ---    \\
    \bottomrule
  \end{tabular}
\end{table}

\end{document}
